Working memory retains information for use under task rules. If memory states have the same content and similar representations, can we use them interchangeably in subsequent computation? Six independently initialized recurrent networks were trained on a permutation-based memory task. Each network received a permutation of 1 through 7, processed operation cues separated by delays, and answered queries about the resulting permutation. The operations included inversion and complementation (replacing each position label j with $8-j$). A preceding context cue determined whether A meant inversion and B complementation, or vice versa. Exact permutation relations allowed different initial inputs and histories to reach identical current permutations. Within each network, we replaced the memory component of the receiver's hidden state with the donor's while retaining the receiver's control component, subsequent inputs, and readout. Exchanges within the same rule context introduced new errors on 1.7\% of all queries, compared with 53.4\% across contexts. Under identical subsequent inputs, answers after exchange matched the donor's own answers on $>99.99$\% of queries, including nearly all donor errors. Some paired control states were identical. Translating subsequent cues into their donor-context equivalents increased cross-context accuracy from 42.1\% to 97.9\% on held-out permutations; the same translation without exchange reduced original-task accuracy from 97.9\% to 42.1\%. Linear centered kernel alignment (CKA) distributions largely overlapped between within-context and cross-context pairs (medians approximately 0.892 and 0.889). Adding CKA to a baseline accounting for task conditions, native performance, content decoding, and input/control-state distances did not consistently improve loss predictions for held-out networks and condition pairs. These results suggest that content-matched memory states can carry different information about how subsequent cues are interpreted. In these networks, matching content and high CKA did not guarantee functional interchangeability, motivating causal tests when interpreting neural population representations.
